% Spiegelpaar zu content/frontmatter/abstract.tex — inhaltlich identisch
% halten (destilliert aus Kap. 5.1 und den RQ-Boxen in 4.7).
Cyber Threat Intelligence (CTI) verlangt von Analysten, Bedrohungs\-informationen
aus heterogenen Quellen unter Zeitdruck zu korrelieren. Large Language
Models (LLMs) k\"onnen diese Arbeit unterst\"utzen, scheitern im direkten
Einsatz jedoch an drei Grenzen: veraltetem Modellwissen, Halluzinationen
und fehlender Provenienz der Aussagen.

Diese Masterarbeit untersucht, wie Retrieval-Augmented Generation (RAG)
die Qualit\"at, Verl\"asslichkeit und operative N\"utzlichkeit von
CTI-Analyse und Threat-Hunting-Workflows verbessern kann. Dazu wird ein
Open-Source-Prototyp konzipiert, implementiert und evaluiert, der vier
kuratierte CTI-Quellen (NVD, CISA KEV, CISA-Advisories, MISP; 84.982
Dokumente) in einen provenienz-erhaltenden Index \"uberf\"uhrt. Ein
hybrides Retrieval kombiniert BM25 und dichte Vektorsuche mit
Cross-Encoder-Reranking. Die Antwortgenerierung ist geschichtet:
Verifizierbare Fakten werden deterministisch aus Passagen-Metadaten
gerendert, ein lokal betriebenes 8B-Modell liefert ausschlie\ss lich das
Narrativ, mit erzwungenen, validierten Zitaten und expliziter
Enthaltung ohne Evidenz.

Die Evaluation vergleicht sechs Konfigurationen -- von der
retrieval-freien Baseline bis zur gehosteten Ablation -- auf einem
eingefrorenen Index-Stand mit 28 Queries \"uber vier
Analyst-Task-Typen und drei L\"aufen je Konfiguration. Bewertet wird
mit RAGAS, Field Coverage und einem vierdimensionalen
Analysten-Rubric samt szenariobasierter Analyse. Zentrales Ergebnis:
Die Antwortkorrektheit bleibt auf Baseline-Niveau, die
Nachvollziehbarkeit (Traceability) steigt jedoch von 1,23 auf 2,56 --
Retrieval liefert auf diesem Query-Set in erster Linie
Verifizierbarkeit, nicht h\"ohere Korrektheit.

Der wissenschaftliche Beitrag umfasst Workflow-\"ubergreifende
empirische Evidenz f\"ur RAG in der CTI, den Befund, dass generische
RAG-Metriken und analystenorientierte Ma\ss e auseinanderfallen, sowie
ein Referenzmodell mit acht empirisch begr\"undeten Designprinzipien.
Der praktische Beitrag ist ein reproduzierbares, lokal betreibbares
System samt quantifiziertem Local-versus-Hosted-Trade-off.
